The automation of software development and delivery processes has become essential for organizations seeking to reduce delivery cycles without compromising quality, stability, and security. However, it is observed that companies face difficulties in adopting these optimization principles in the development and delivery of pipelines within their processes. Therefore, this work developed an architecture for CI/CD pipeline automation based on Kubernetes, considering performance, evaluation, and security mechanisms under Zero Trust principles. The research adopts an experimental approach, grounded in a literature review and the application of DORA metrics and Zero Trust for performance analysis. The experiment includes the implementation of a CI/CD architecture on Kubernetes with ephemeral agents and different computational infrastructure configurations and security mechanisms.
The results indicate that the proposed architecture contributes to greater scalability, isolation of executions, and a reduction in computational costs of up to 90% through the segregation of workloads into Spot instances. From a security perspective, the experiments evidenced the mitigation of 100% of compromise propagation threats through the automatic isolation of non-compliant services via SPIRE. It was also observed that validation mechanisms (such as Cosign and Rekor) add low operational impact, adding an average latency of only 6 seconds to the attestation time, while significantly increasing traceability and trust in deployed artifacts. Regarding DORA metrics, the solution showed performance compatible with the "Elite" maturity level, demonstrating the potential to reduce Lead Time to less than one hour and maintain change failure rates between 0% and 15%. It is concluded that the integration of DevOps, Kubernetes, DORA metrics, and Zero Trust constitutes a viable approach for building modern continuous delivery platforms.

\textbf{Keywords}: DevOps. CI/CD. Kubernetes. DORA Metrics. Zero Trust.